% Experimental results (V5 refocused revision).
% V5 directive: FULL graphical + tabular treatment for Phase B1 and Phase C2
% ONLY. Every other phase (A, C, D, B2, B2-nr, multi-seed, calibration,
% throughput, confusion) is collapsed to a short qualitative paragraph with no
% tables and no figures. For Phase C2 the degradation axes are compared on the
% measured image-quality scales (PSNR and SSIM) -- never on the L1-L5 level
% index, which is a predetermined parameter rather than a quality measure.
\section{Experimental Results}
\label{sec:tests}

This section reports the measured performance against every relevant metric.
Every value comes from a frozen training run recorded under
\texttt{runs/final/} and aggregated by \texttt{artifacts/Final\_Exp.json}
(schema version 3). No values are synthesized or interpolated. In line with
the focus of this revision, the two phases that carry the headline science ---
Phase B1 (combined degradation) and Phase C2 (single-axis isolation under the
THz protocol) --- are reported in full, with tables, figures, and a
model-by-model reading. The remaining phases are summarized qualitatively in
Section~\ref{sec:tests:brief}; their complete per-cell numbers live in the
aggregated artifacts referenced in the Appendix.

\subsection{Test rig}
\label{sec:tests:rig}
The campaign ran on a single workstation with the configuration in
Table~\ref{tab:rig}. Determinism is guarded by
\texttt{pl.seed\_everything(42, workers=True)} for the canonical seed and by
the per-sample seed offset for the validation split.

\begin{table}[htbp]
\centering
\caption{Test rig and software stack.}
\label{tab:rig}
\setlength{\tabcolsep}{4pt}
\renewcommand{\arraystretch}{1.1}
\begin{tabular}{ll}
\toprule
Component & Value \\
\midrule
GPU                & NVIDIA RTX 5070 (Blackwell, sm\_120, 12 GB) \\
Precision          & \texttt{bf16-mixed} \\
CPU dataloader workers & 0 (Windows) \\
Framework          & PyTorch 2.11 [17] + Lightning 2.x [16] \\
Image quality      & \texttt{torchmetrics} PSNR / SSIM, \texttt{data\_range = 1} \\
Determinism gate   & \texttt{src/tests/test\_degradation\_determinism.py} \\
\bottomrule
\end{tabular}
\end{table}

\subsection{Degradation examples}
\label{sec:tests:samples}
Figure~\ref{fig:samples} grounds the rest of the section visually. It shows
six fixed validation examples per dataset under the clean input, the combined
Phase B1 pipeline at the moderate level, and each of the three Phase C2 axes
(resolution, blur, salt-and-pepper) isolated under the THz protocol
(saturation and additive noise disabled).

\begin{figure}[htbp]
  \centering
  \includegraphics[width=0.49\textwidth]{\figpath{samples/grid_phase_THz_L3_cifar10.png}}\hfill
  \includegraphics[width=0.49\textwidth]{\figpath{samples/grid_phase_THz_L3_mnist.png}}
  \caption{Visual examples of the degradation pipelines. Left: CIFAR-10. Right: MNIST. Rows: clean input; combined Phase B1 at L3; and the three Phase C2 axes (resolution, blur, salt-and-pepper) isolated at L3 under the THz protocol.}
  \label{fig:samples}
\end{figure}

\subsection{The remaining phases in brief}
\label{sec:tests:brief}
The phases outside the B1 / C2 focus are summarized here qualitatively; they
provide context and bounds for the two detailed analyses that follow.

\sectheading{Clean baselines (Phase A).} All three backbones classify the
undegraded validation sets at near-ceiling accuracy, establishing the upper
bound against which every degraded cell is measured.

\sectheading{Full single-axis isolation (Phase C).} Sweeping all five
degradation axes one at a time confirms that spatial resolution is the dominant
driver of accuracy loss, whereas desaturation and salt-and-pepper noise are
comparatively benign in isolation. This ordering motivated the THz-relevant
three-axis subset that is analysed in depth as Phase C2.

\sectheading{Regularization recovery (Phase D).} Layering dropout, mix-up and
cut-mix on top of the combined-degradation pipeline recovers only a small
fraction of the accuracy lost to degradation. The collapse is therefore an
information bottleneck, not an over-fitting artefact, and cannot be closed by
regularization alone.

\sectheading{THz protocol simplification (Phase B2 / B2-nr).} Reducing the
pipeline to the THz-relevant operators --- removing the colour-desaturation and
additive-noise axes --- yields a modest, consistent accuracy margin over the
full combined pipeline. Decomposing that margin shows it is driven mostly by
the protocol simplification itself rather than by regularization.

\sectheading{Multi-seed audit.} Re-running the headline cells under additional
random seeds leaves the reported accuracies essentially unchanged: run-to-run
variance is small relative to every effect discussed below, so the comparisons
are robust to seed choice.

\sectheading{Calibration.} The networks grow over-confident as the spatial
information collapses --- predicted confidence increasingly exceeds empirical
accuracy at the most severe levels. Regularization and protocol simplification
reduce but do not eliminate this miscalibration.

\sectheading{Inference throughput.} Inference latency meets the real-time
budget for the convolutional backbones; the transformer backbone is somewhat
slower because its attention path is not graph-compiled on the current stack.

\sectheading{Confusion structure.} At the most severe level the residual errors
concentrate among visually similar classes, indicating the networks retain
coarse category structure even after fine detail is destroyed.

\subsection{Phase B1: combined degradation}
\label{sec:tests:phase_b1}
Phase B1 applies all five degradation axes simultaneously at a common level.
Table~\ref{tab:phase_b} reports the full $3 \times 2 \times 5$ accuracy block,
and Figure~\ref{fig:b1_curves} plots each model's accuracy against the combined
level, with the Phase A clean baseline drawn as a dotted line per model. Here
the level index is a legitimate $x$-axis: Phase B1 sweeps a single fixed
combined recipe, so each level denotes one well-defined operating point.

\begin{table}[htbp]
\centering
\caption{Phase B (combined degradation) validation accuracy (\%). All five axes active at the same level L. Cells with $\pm$ values are multi-seed audited (seeds $\{42, 43, 44\}$).}
\label{tab:phase_b}
\setlength{\tabcolsep}{4pt}
\renewcommand{\arraystretch}{1.1}
\begin{tabular}{llrrrrr}
\toprule
Model & Dataset & L1 & L2 & L3 & L4 & L5 \\
\midrule
ResNet50 & CIFAR-10 & 76.46 & 53.94 & 36.31 $\pm$ 0.83 & 26.84 & 19.06 \\
ResNet50 & MNIST & 99.18 & 96.36 & 79.17 $\pm$ 0.80 & 56.44 & 28.00 \\
DenseNet121 & CIFAR-10 & 78.48 & 53.68 & 38.54 $\pm$ 0.78 & 28.52 & 19.58 \\
DenseNet121 & MNIST & 98.74 & 96.54 & 80.79 $\pm$ 0.84 & 58.08 & 27.76 \\
TransNeXt-tiny & CIFAR-10 & 87.58 & 64.22 & 39.48 $\pm$ 0.62 & 28.50 & 20.10 \\
TransNeXt-tiny & MNIST & 98.96 & 96.00 & 80.05 $\pm$ 0.11 & 57.24 & 27.48 \\
\bottomrule
\end{tabular}
\end{table}


\begin{figure}[htbp]
  \centering
  \includegraphics[width=0.49\textwidth]{\figpath{curves/b1_acc_vs_level_cifar10.png}}\hfill
  \includegraphics[width=0.49\textwidth]{\figpath{curves/b1_acc_vs_level_mnist.png}}
  \caption{Phase B1 validation accuracy versus combined degradation level. Solid curves: the three backbones; dotted horizontal lines: the per-model Phase A clean baseline. Left: CIFAR-10. Right: MNIST.}
  \label{fig:b1_curves}
\end{figure}

On CIFAR-10 accuracy declines monotonically from the mild level (L1) to the
extreme level (L5). TransNeXt-tiny leads at the mild and moderate end
($87.58\,\%$ at L1, against a CNN mean near $77\,\%$), but the three
architectures converge into a tight $[19,21]\,\%$ band at L5, where the spatial
information has effectively been removed. MNIST is far more resilient: all
three backbones stay above $98\,\%$ at L1 and only fall to roughly
$27$--$28\,\%$ at L5. The model-by-model reading of this convergence is given
in Section~\ref{sec:crossmodel}.

\subsection{Phase C2: single-axis under the THz protocol}
\label{sec:tests:phase_c2}
Phase C2 isolates each of the three THz-relevant degradation axes ---
resolution, blur, and salt-and-pepper --- under the simplified protocol
(saturation and additive noise disabled), so each cell varies exactly one
operator. Table~\ref{tab:phase_c2} reports the accuracy block.

\begin{table}[htbp]
\centering
\caption{Phase C2 single-axis isolation under the THz protocol (saturation $=0$, noise $=0$, T3 regularization) validation accuracy (\%).}
\label{tab:phase_c2}
\setlength{\tabcolsep}{4pt}
\renewcommand{\arraystretch}{1.1}
\begin{tabular}{lllrrrrr}
\toprule
Model & Dataset & Axis & L1 & L2 & L3 & L4 & L5 \\
\midrule
ResNet50 & CIFAR-10 & Resolution & 80.18 & 66.36 & 51.32 & 42.96 & 30.16 \\
ResNet50 & CIFAR-10 & Blur & 90.62 & 89.58 & 83.36 & 71.68 & 58.76 \\
ResNet50 & CIFAR-10 & S\&P & 86.24 & 85.18 & 83.48 & 77.86 & 74.68 \\
ResNet50 & MNIST & Resolution & 99.06 & 98.64 & 94.12 & 81.24 & 45.54 \\
ResNet50 & MNIST & Blur & 99.24 & 99.26 & 99.32 & 99.32 & 98.86 \\
ResNet50 & MNIST & S\&P & 99.14 & 99.16 & 99.18 & 99.00 & 99.00 \\
DenseNet121 & CIFAR-10 & Resolution & 76.82 & 64.04 & 51.34 & 44.44 & 30.80 \\
DenseNet121 & CIFAR-10 & Blur & 88.24 & 85.66 & 81.66 & 74.90 & 64.60 \\
DenseNet121 & CIFAR-10 & S\&P & 86.38 & 83.26 & 80.54 & 79.30 & 77.78 \\
DenseNet121 & MNIST & Resolution & 99.38 & 98.64 & 94.30 & 82.28 & 45.02 \\
DenseNet121 & MNIST & Blur & 99.42 & 99.34 & 99.22 & 99.38 & 99.16 \\
DenseNet121 & MNIST & S\&P & 99.16 & 99.20 & 99.20 & 99.08 & 99.06 \\
TransNeXt-tiny & CIFAR-10 & Resolution & 89.34 & 76.08 & 58.24 & 46.58 & 30.54 \\
TransNeXt-tiny & CIFAR-10 & Blur & 96.52 & 95.96 & 93.00 & 88.02 & 71.04 \\
TransNeXt-tiny & CIFAR-10 & S\&P & 95.20 & 94.78 & 93.52 & 91.56 & 89.58 \\
TransNeXt-tiny & MNIST & Resolution & 99.32 & 98.76 & 93.94 & 81.74 & 45.60 \\
TransNeXt-tiny & MNIST & Blur & 99.38 & 99.38 & 99.46 & 99.48 & 99.06 \\
TransNeXt-tiny & MNIST & S\&P & 99.28 & 99.14 & 99.20 & 99.08 & 99.08 \\
\bottomrule
\end{tabular}
\end{table}


A given nominal level does not impose the same image quality on different
axes: L3 blur and L3 resolution are not equally destructive. Comparing the
axes on the L1--L5 index would therefore be meaningless --- the level is a
predetermined parameter, not a measure of how much the image was actually
degraded. We instead plot accuracy against the \emph{measured} image quality,
PSNR and SSIM, computed against the clean upsampled reference.

Figure~\ref{fig:c2_type} compares the three axes on a common PSNR and SSIM
scale, with one curve per axis averaged over the three models. At matched image
quality, resolution degradation is by far the most harmful: its curve sits well
below blur and salt-and-pepper at every PSNR and SSIM value, while
salt-and-pepper is nearly flat. The Spearman rank correlation confirms that
resolution carries the strongest accuracy--quality association
($\rho = +0.98$ against both metrics on CIFAR-10).

\begin{figure}[htbp]
  \centering
  \begin{tabular}{cc}
    \includegraphics[width=0.46\textwidth]{\figpath{curves/c2_type_psnr_cifar10.png}} &
    \includegraphics[width=0.46\textwidth]{\figpath{curves/c2_type_ssim_cifar10.png}} \\
    {\small (a) CIFAR-10, vs PSNR} & {\small (b) CIFAR-10, vs SSIM} \\[0.3em]
    \includegraphics[width=0.46\textwidth]{\figpath{curves/c2_type_psnr_mnist.png}} &
    \includegraphics[width=0.46\textwidth]{\figpath{curves/c2_type_ssim_mnist.png}} \\
    {\small (c) MNIST, vs PSNR} & {\small (d) MNIST, vs SSIM} \\
  \end{tabular}
  \caption{Phase C2 degradation-type comparison. Validation accuracy against measured image quality (PSNR, left column; SSIM, right column), one curve per degradation axis, averaged over the three backbones. At equal image quality the resolution axis is consistently the most destructive and salt-and-pepper the least. The level index is deliberately not used as the $x$-axis.}
  \label{fig:c2_type}
\end{figure}

Figure~\ref{fig:c2_model} repeats the analysis per axis with one curve per
model, against PSNR; Figure~\ref{fig:c2_model_ssim} gives the matching SSIM
view. TransNeXt-tiny is the most robust on every axis, the two CNNs essentially
overlap, and the model ordering is preserved across the entire quality range.
On MNIST the digit shapes tolerate blur and salt-and-pepper almost completely,
so only the resolution axis drives a meaningful accuracy loss.

\begin{figure}[htbp]
  \centering
  \includegraphics[width=\textwidth]{\figpath{curves/c2_model_psnr_cifar10.png}}\\[0.4em]
  \includegraphics[width=\textwidth]{\figpath{curves/c2_model_psnr_mnist.png}}
  \caption{Phase C2 model comparison versus PSNR. One panel per axis (resolution, blur, salt-and-pepper); one curve per backbone. Top: CIFAR-10. Bottom: MNIST.}
  \label{fig:c2_model}
\end{figure}

\begin{figure}[htbp]
  \centering
  \includegraphics[width=\textwidth]{\figpath{curves/c2_model_ssim_cifar10.png}}\\[0.4em]
  \includegraphics[width=\textwidth]{\figpath{curves/c2_model_ssim_mnist.png}}
  \caption{Phase C2 model comparison versus SSIM. One panel per axis; one curve per backbone. Top: CIFAR-10. Bottom: MNIST.}
  \label{fig:c2_model_ssim}
\end{figure}
